\subsubsection{Adjacent-zero shielding under backward heat flow}
\label{subsec:zero-shielding}
A pointwise Laguerre-source comparison is sufficient to
preserve positivity, but the preceding collision analysis
shows why the required coefficient can diverge. There is
a complementary route that does not divide by Laguerre
expressions: follow the \emph{distances between consecutive
real zeros} while they exist.

The zero-motion equation for the de Bruijn--Newman flow
is classical, going back to Csordas, Smith and Varga
and playing a central role in Rodgers and Tao
\cite{RodgersTao}. The elementary subtraction below
isolates an explicit interaction between a given gap
and all the other zeros. We record it to clarify
the quantitative continuation problem, not as a
new zero-motion theorem.

\begin{proposition}[Exact neighboring-gap identity]
\label{prop:gapshield}
On an interval of heat parameters where all zeros of
$H_t$ are real and simple, take two consecutive zeros
$a(t)<b(t)$, and put $g(t)=b(t)-a(t)$. Write
$s=\frac12-t$ and
\[
\mathcal S_{a,b}(t):=
\sum_{\rho(t)\ne a(t),b(t)}
\frac{1}{(a(t)-\rho(t))(b(t)-\rho(t))}.
\]
The sum converges absolutely for each such pair,
and is strictly positive. With $G(s)=g(\frac12-s)^2$,
\begin{equation}
\boxed{G'(s)=-8+4G(s)\,\mathcal S_{a,b}(\tfrac12-s).}
\label{eq:shieldgap}
\end{equation}
Consequently, for any $0\le \sigma$ inside the
real-simple-zero regime,
\begin{equation}
\boxed{
G(\sigma)=G(0)-8\sigma+
4\int_0^\sigma G(v)\mathcal S_{a,b}(\tfrac12-v)\,dv.
}
\label{eq:shieldbudget}
\end{equation}
In particular, a nonnegative endpoint gap necessarily
requires
\[
4\int_0^\sigma G(v)\mathcal S_{a,b}(\tfrac12-v)\,dv
\ge 8\sigma-G(0).
\]
\end{proposition}
\begin{proof}
The classical zero-motion equation is
\[
\dot a(t)=2\,\mathrm{PV}\sum_{\rho\ne a}
\frac1{a-\rho},\qquad
\dot b(t)=2\,\mathrm{PV}\sum_{\rho\ne b}
\frac1{b-\rho}.
\]
The two terms contributed by the opposite endpoints
give $4/g$ to $\dot b-\dot a$. For any other real zero,
\[
\frac1{b-\rho}-\frac1{a-\rho}
=-\frac{g}{(a-\rho)(b-\rho)}.
\]
The principal-value issue disappears in the
\emph{difference}: the terms converge absolutely
because the zero counting function has the usual
$O(T\log T)$ growth. Thus
$\dot g=4/g-2g\mathcal S_{a,b}$.
Changing to $s=\frac12-t$ and differentiating
$G=g^2$ proves \eqref{eq:shieldgap}.
Integration gives \eqref{eq:shieldbudget}.
For consecutive real roots every other $\rho$
lies outside $[a,b]$, so each denominator product
is positive.
\end{proof}

\begin{proposition}[Exact cancellation in a regular lattice]
\label{prop:latticeshield}
For an ideal arithmetic progression of zeros
$\rho_k=kg$ with consecutive endpoints $0,g$,
\[
\boxed{
g^2\mathcal S_{0,g}
=2\sum_{m=1}^\infty\frac1{m(m+1)}=2.}
\]
Therefore the right side of \eqref{eq:shieldgap}
vanishes identically. If one retains only
the closest $R$ roots on each side,
\begin{equation}
g^2\mathcal S_R=\frac{2R}{R+1},
\qquad
-8+4g^2\mathcal S_R=-\frac8{R+1}.
\label{eq:truncatedshield}
\end{equation}
The deficit $8/(R+1)$ comes entirely from
omitting the remaining zeros.
\end{proposition}
\begin{proof}
The contribution from the $m$th exterior
zero on either side is
$1/(g^2m(m+1))$.
Use $1/(m(m+1))=1/m-1/(m+1)$.
The infinite sum is $1$ and
the partial sum is $R/(R+1)$.
\end{proof}

The regular lattice is an exact heat-flow control:
$e^{\kappa^2t}\sin(\kappa z)$ solves
$\partial_tH=-\partial_z^2H$
and has a stationary arithmetic zero lattice.
It is not a positive theta Fourier transform.
The cancellation in Proposition~\ref{prop:latticeshield}
is thus a model identity, not a statement that
the Riemann zeros form an arithmetic progression.

\begin{theorem}[Conditional finite-window shielding estimate]
\label{thm:nearshieldsurvival}
Suppose the relevant zeros are real and simple
on a backward-time interval $0\le s<S$.
Fix a consecutive pair with gap $g(s)>0$.
Let $R\ge1$ and $0\le\varepsilon<1$.
Assume the $m$th root left of $a(s)$
and the $m$th root right of $b(s)$,
for each $1\le m\le R$, are within
$\varepsilon g(s)$ of their ideal
positions $a(s)-mg(s)$ and $b(s)+mg(s)$,
respectively, throughout that interval.
Then
\begin{equation}
\frac{2R}{(R+1)(1+\varepsilon)^2}
\le g(s)^2\mathcal S_R(s)
\le\frac{2R}{(R+1)(1-\varepsilon)^2},
\label{eq:latticeeps}
\end{equation}
and the full squared gap satisfies
\begin{equation}
\boxed{
G'(s)\ge-\frac8{R+1}-16\varepsilon,\qquad
G(s)\ge G(0)-s
\left(\frac8{R+1}+16\varepsilon\right).
}
\label{eq:gapstability}
\end{equation}
In particular, throughout any such interval of
length at most $\tfrac12$, a sufficient
\emph{conditional} positive-margin criterion is
\begin{equation}
\boxed{g(0)^2>\frac4{R+1}+8\varepsilon.}
\label{eq:shieldcriterion}
\end{equation}
One explicit choice satisfying this numerical margin is
$R\ge16/g(0)^2$ and
$\varepsilon\le g(0)^2/32$, provided $\varepsilon<1$.
\end{theorem}
\begin{proof}
For the left $m$th exterior root, the two
distances to the endpoints lie between
$(m-\varepsilon)g$ and $(m+\varepsilon)g$,
and between $(m+1-\varepsilon)g$ and
$(m+1+\varepsilon)g$.
The same inequalities hold in absolute value
for the right exterior root. Thus each
positive summand of $g^2\mathcal S_R$
lies between
\[
\frac{1}{(1+\varepsilon)^2m(m+1)}
\quad\text{and}\quad
\frac{1}{(1-\varepsilon)^2m(m+1)}.
\]
Summing and using the telescoping identity
gives \eqref{eq:latticeeps}.
All omitted exterior roots contribute
positively, hence
$G'\ge-8+8R/((R+1)(1+\varepsilon)^2)$.
The elementary bound
$1-(1+\varepsilon)^{-2}\le2\varepsilon$
gives the first inequality in
\eqref{eq:gapstability}; integration gives
the second and \eqref{eq:shieldcriterion}.
For the displayed parameter choice,
$4/(R+1)<g(0)^2/4$ and
$8\varepsilon\le g(0)^2/4$.
\end{proof}

\paragraph{The missing theta-specific estimate.}
The hypotheses of Theorem~\ref{thm:nearshieldsurvival}
have \emph{not} been established for the Riemann
theta zeros as $t$ decreases from $\tfrac12$ to $0$.
In fact, the required precision is severe.
At height $X$ a typical gap is on the scale
$g\asymp1/\log X$, so the displayed sufficient
condition requires $R$ of order at least
$(\log X)^2$ and relative lattice deviations
$\varepsilon$ of order at most $(\log X)^{-2}$.
Such pointwise rigidity should not be confused
with average spacing laws or pair-correlation
estimates. It may be far stronger than necessary.

The more flexible exact target is the
\emph{integrated shielding defect}
\[
\mathcal D_{a,b}(s)
:=2-G(s)\mathcal S_{a,b}(\tfrac12-s),
\qquad G'(s)=-4\mathcal D_{a,b}(s).
\]
To rule out a finite collision of this pair
before backward time $S$, it would suffice,
while the real-zero dynamics are valid, to
establish
\begin{equation}
\boxed{
4\int_0^\sigma
[\mathcal D_{a,b}(v)]_+\,dv<G(0)
\qquad(0\le\sigma\le S),
}
\label{eq:integratedshieldu}
\end{equation}
with a margin uniform across pairs and heights.
A global RH deduction would further require
control of possible loss of real-rootedness
through zeros escaping to unbounded height.
Neither estimate is known here.
The equation of motion and regular-lattice
equilibrium are classical aspects of the
de Bruijn--Newman programme; the finite-window
bounds are elementary conditional consequences,
not a claim of a proof of RH.

\paragraph{Exact verification.}
The independent script
\path{certificates/xi_heat_zero_shielding_exact.py}
checks the telescoping sums to $R=500$,
the heat-root ODE at 14 adjacent pairs
of rational-root polynomials, and rational
near-lattice inequalities for adversarial
displacements. The Lean file
\path{XiKernel/HeatZeroShielding.lean}
checks the algebraic gap, shielding-budget
and positive-margin implications only.
Neither the genuine infinite zero-motion
equation nor the necessary all-height
theta-specific shielding estimate is
currently formalized or proved end to end.
